
\documentclass{article}
\usepackage{colortbl}
\usepackage{makecell}
\usepackage{multirow}
\usepackage{supertabular}

\begin{document}

\newcounter{utterance}

\centering \large Interaction Transcript for game `wordle\_withclue', experiment `medium\_frequency\_words\_no\_clue\_no\_critic', episode 1 with reka{-}flash{-}3.1{-}t0.0.
\vspace{24pt}

{ \footnotesize  \setcounter{utterance}{1}
\setlength{\tabcolsep}{0pt}
\begin{supertabular}{c@{$\;$}|p{.15\linewidth}@{}p{.15\linewidth}p{.15\linewidth}p{.15\linewidth}p{.15\linewidth}p{.15\linewidth}}
    \# & \multicolumn{2}{c}{Player} && \multicolumn{2}{c}{Game Master} \\
    \hline

    \theutterance \stepcounter{utterance}  
    & & & \multicolumn{4}{p{0.6\linewidth}}{
        \cellcolor[rgb]{0.9,0.9,0.9}{
            \makecell[{{p{\linewidth}}}]{
                \texttt{\tiny{[P1$\langle$GM]}}
                \texttt{You are a language wizard who likes to guess words by using the given rules.} \\
\\ 
\texttt{Welcome to Wordle! You have six attempts to guess the target word, a valid English word of five lowercase letters (a{-}z). Please use the tags "explanation:" and "guess:" to provide a concise explanation for each guess.} \\
\\ 
\texttt{For instance, if your guess is "apple", your response should be} \\
\texttt{explanation: this is a common five{-}letter English word, and I am starting my guess with this word.} \\
\texttt{guess: apple} \\
\\ 
\texttt{After each guess, your answer will be validated, and you will receive feedback indicating which letters are correct (green), which letters are correct but in the wrong position (yellow), and which letters are incorrect (red). This feedback can be useful in determining which letters to include or exclude in your next guess.} \\
\\ 
\texttt{For example, the feedback for "apple" might be:} \\
\texttt{guess\_feedback: a<yellow> p<yellow> p<green> l<yellow> e<red>} \\
\\ 
\texttt{The explanation should contain details about how the guess\_feedback is used to arrive at a new guess.} \\
\\ 
\texttt{Do not generate any other text.} \\
\\ 
\texttt{Let's begin with your first guess.} \\
\\ 
\texttt{clue: logbook} \\
            }
        }
    }
     \\ \\

    \theutterance \stepcounter{utterance}  
    & \multicolumn{4}{p{0.6\linewidth}}{
        \cellcolor[rgb]{0.561,0.737,0.561}{
            \makecell[{{p{\linewidth}}}]{
                \texttt{\tiny{[P1$\rangle$GM]}}
                \texttt{<reasoning>} \\
\texttt{Let's tackle this Wordle. The first clue given is "logbook". I need to figure out how to use that. Maybe "logbook" is a hint about the target word? Or perhaps it's a clue in another way. Let me think.} \\
\\ 
\texttt{Wait, maybe "logbook" is a red herring, and the actual target is a five{-}letter word. But the user mentioned that the target is a five{-}letter word. So maybe "logbook" is a clue to help me start. Let's break down "logbook". It has letters L, O, G, B, O, O, K. Maybe some letters are repeated here. The letters present are L, O, G, B, K. Since the target is five letters, maybe some of these letters are in the target. But how?} \\
\\ 
\texttt{Alternatively, maybe "logbook" is a clue that the target word is related to it. For example, "logbook" is a record, so maybe the target is a word like "record"? But "record" is six letters. Hmm. Not helpful.} \\
\\ 
\texttt{Wait, maybe the clue is that the target word is a part of "logbook". Let's see: logbook has letters L, O, G, B, O, O, K. The target is five letters. Maybe the target is "globe"? But that's five letters. Let me check if "globe" is a valid word. Yes, it is. But why would "logbook" hint at "globe"? Maybe because "globe" is part of "logbook"? Not sure. Alternatively, maybe the letters in "logbook" are clues. Let's list the unique letters: L, O, G, B, K. So the target word must include some of these letters. But since the target is five letters, maybe it's a combination of these letters. But how?} \\
\\ 
\texttt{Alternatively, maybe the clue is that the target word is a synonym or related term. "Logbook" is a book of logs, maybe related to navigation or aviation. Words like "chart", "route", "plane", "pilot", "naval", etc. But those might not be five letters. Let's see: "chart" is five letters. "Route" is five. "Plane" is five. "Pilot" is six. "Naval" is five. But how does "log} \\
            }
        }
    }
    & & \\ \\

    \theutterance \stepcounter{utterance}  
    & & & \multicolumn{2}{p{0.3\linewidth}}{
        \cellcolor[rgb]{0.9,0.9,0.9}{
            \makecell[{{p{\linewidth}}}]{
                \texttt{\tiny{[GM$|$GM]}}
                \texttt{Error: The response should always start with the keyword 'explanation:'} \\
            }
        }
    }
    & & \\ \\

    \theutterance \stepcounter{utterance}  
    & & & \multicolumn{2}{p{0.3\linewidth}}{
        \cellcolor[rgb]{0.9,0.9,0.9}{
            \makecell[{{p{\linewidth}}}]{
                \texttt{\tiny{[GM$|$GM]}}
                \texttt{game\_result = ABORT} \\
            }
        }
    }
    & & \\ \\

\end{supertabular}
}

\end{document}
